\documentclass[10pt,letterpaper]{article}
\usepackage{cornell-notes}

\title{Clause 16.03.02.04: Detailed Specifications}
\author{}
\date{}
\setDocTitle{Clause 16.03.02.04: Detailed Specifications}
\setDocSubtitle{C++ 2024}
\setDocOwner{C++ 2024 Cornell Notes}
\setDocDate{\today}
\setCornellCollection{C++ 2024 Cornell Notes}
\setCornellUnitType{Clause}
\setCornellUnitNumber{16.03.02.04}
\setCornellUnitTitle{Detailed Specifications}
\hypersetup{pdftitle={Clause 16.03.02.04: Detailed Specifications},pdfsubject={Cornell notes},pdfauthor={}}

\begin{document}
\maketitle
\begin{CornellOverviewBox}{Overview}
\textbf{Classification:} Standard library - Normative\\[2pt]
\textbf{Learning objective:} Understand the rules, terminology, and semantic consequences associated with Detailed specifications.
\end{CornellOverviewBox}\begin{CornellSummaryBox}{Summary}
{\sffamily\bfseries\color{Accent} Summary}\par\vspace{3pt}Section 16.3.2.4 focuses on Detailed specifications. Apply the specified interface together with its mandates, effects, result conditions, exception behavior, and complexity constraints. When applying it, preserve the distinction between mandatory requirements, recommendations, permissions, and behavior deliberately left undefined, unspecified, or implementation-defined.
\end{CornellSummaryBox}

\end{document}
